% TODO checklist:
% - [x] Inspected pyproject.toml for Python version and package config
% - [x] Inspected requirements.txt for all runtime dependencies
% - [x] Inspected docker-compose.yml for service definitions
% - [x] Inspected ui_agent/package.json for frontend stack
% - [x] Inspected ui_control_panel/ for Streamlit dependencies
% - [x] Inspected src/config.py for LLM provider configuration
% - [x] Referenced README.md and Q&A.md for architectural overview

\section{Technology Stack}
\label{sec:tech-stack}

The Cineca Agentic Platform is built on a modern, production-ready technology stack carefully selected to balance performance, maintainability, and enterprise requirements. This section provides a comprehensive overview of the technologies employed across all platform layers.

\subsection{Backend Core Technologies}

The platform's backend is implemented primarily in Python, leveraging the asynchronous capabilities of modern Python frameworks for high-throughput API handling.

\subsubsection{Python Runtime and Framework}

\begin{table}[ht]
\centering
\caption{Core backend technologies and versions.}
\label{tab:backend-stack}
\small
\begin{tabular}{lll}
\hline
\textbf{Technology} & \textbf{Version} & \textbf{Purpose} \\
\hline
Python & $\geq$ 3.10 & Runtime environment \\
FastAPI & $\geq$ 0.111.0 & Web framework and API layer \\
Uvicorn & $\geq$ 0.30.0 & ASGI server with standard extras \\
Pydantic & $\geq$ 2.2.1 & Data validation and settings management \\
SQLAlchemy & $\geq$ 2.0.30 & ORM for PostgreSQL \\
Alembic & $\geq$ 1.13.0 & Database migrations \\
\hline
\end{tabular}
\end{table}

FastAPI was selected as the primary web framework due to its native support for asynchronous request handling, automatic OpenAPI documentation generation, and Pydantic-based request/response validation. The framework's dependency injection system enables clean separation of concerns and facilitates testing.

\begin{lstlisting}[language=Python, caption={FastAPI application factory pattern from \texttt{src/app.py}.}]
def create_app() -> FastAPI:
    """Create the FastAPI application with all middleware and routers."""
    # Initialize structured logging
    from src.logging_setup import setup_logging
    setup_logging(level=settings.LOG_LEVEL)
    
    # Install sensitive data masking for logs
    from src.security.secrets import install_log_masking
    install_log_masking()
    
    # Create FastAPI instance
    app = FastAPI(
        title="Cineca Agentic Platform",
        version="0.1.0",
        docs_url=None,  # Custom versioned docs
        redoc_url=None,
        openapi_url="/openapi.json",
    )
    return app
\end{lstlisting}

\subsubsection{HTTP Client and Resilience}

For outbound HTTP communication with LLM providers and external services, the platform employs:

\begin{itemize}
    \item \textbf{httpx} ($\geq$ 0.27.0): Modern async-capable HTTP client with connection pooling
    \item \textbf{tenacity} ($\geq$ 8.2.3): Retry library for implementing exponential backoff and circuit breaker patterns
\end{itemize}

\subsection{Data Persistence Layer}

The platform employs a polyglot persistence strategy with three complementary database systems, each optimized for specific workload characteristics.

\subsubsection{PostgreSQL Control Plane}

PostgreSQL serves as the authoritative control plane database, storing all configuration, audit logs, and transactional data.

\begin{table}[ht]
\centering
\caption{PostgreSQL stack components.}
\label{tab:postgres-stack}
\small
\begin{tabular}{lll}
\hline
\textbf{Component} & \textbf{Version} & \textbf{Purpose} \\
\hline
PostgreSQL & 16-alpine & Relational database engine \\
psycopg2-binary & $\geq$ 2.9.9 & PostgreSQL adapter for Python \\
SQLAlchemy & $\geq$ 2.0.30 & ORM and connection pooling \\
Alembic & $\geq$ 1.13.0 & Schema migration management \\
\hline
\end{tabular}
\end{table}

The database configuration supports enterprise deployment requirements through configurable connection pooling:

\begin{lstlisting}[language=Python, caption={PostgreSQL configuration from \texttt{src/config.py}.}]
# PostgreSQL connection settings
DB_HOST: str = Field(default="postgres")
DB_PORT: int = Field(default=5432)
DB_NAME: str = Field(default="cineca_platform")
DB_POOL_SIZE: int = Field(default=10)
DB_POOL_TIMEOUT: int = Field(default=30)
DB_POOL_RECYCLE: int = Field(default=3600)
DB_POOL_PRE_PING: bool = Field(default=True)
\end{lstlisting}

\subsubsection{Redis Data Plane}

Redis provides the low-latency data plane for caching, queuing, and ephemeral state management:

\begin{table}[ht]
\centering
\caption{Redis usage categories in the platform.}
\label{tab:redis-usage}
\small
\begin{tabular}{ll}
\hline
\textbf{Usage Category} & \textbf{Key Pattern} \\
\hline
Job queues & \texttt{jobs:queue:*} \\
Idempotency cache & \texttt{jobs:idempotency:*} \\
Rate limiting & \texttt{ratelimit:*} \\
Session state & \texttt{sessions:*} \\
Default model cache (DMR) & \texttt{dmr:*} \\
JWKS key cache & \texttt{jwks:*} \\
Cancel flags & \texttt{jobs:cancel:*} \\
\hline
\end{tabular}
\end{table}

\subsubsection{Memgraph Graph Database}

Memgraph serves as the knowledge graph backend, storing domain entities and relationships for natural language querying:

\begin{itemize}
    \item \textbf{Memgraph Platform}: Latest version with Lab visualization interface
    \item \textbf{gqlalchemy} ($\geq$ 1.8.0): Python client and ORM for Memgraph
    \item \textbf{Cypher}: Query language for graph traversal and manipulation
\end{itemize}

\subsection{Security and Authentication}

The security stack implements industry-standard authentication and authorization protocols:

\begin{table}[ht]
\centering
\caption{Security technologies and standards.}
\label{tab:security-stack}
\small
\begin{tabular}{lll}
\hline
\textbf{Technology} & \textbf{Version} & \textbf{Purpose} \\
\hline
python-jose & $\geq$ 3.3.0 & JWT encoding/decoding with cryptography \\
passlib & $\geq$ 1.7.4 & Password hashing with bcrypt \\
email-validator & $\geq$ 2.1.1 & Email format validation \\
\hline
\end{tabular}
\end{table}

The platform supports OIDC/JWT-based authentication through configurable identity providers:

\begin{lstlisting}[language=Python, caption={OIDC configuration settings.}]
# OIDC Resource Server Settings
OIDC_ISSUER: str | None = Field(default=None)
OIDC_AUDIENCE: str | None = Field(default=None)
OIDC_JWKS_URL: str | None = Field(default=None)
OIDC_TIMEOUT_S: int = Field(default=5)
\end{lstlisting}

\subsection{Observability Stack}

The platform implements comprehensive observability through three pillars: metrics, logging, and tracing.

\begin{table}[ht]
\centering
\caption{Observability technologies.}
\label{tab:observability-stack}
\small
\begin{tabular}{lll}
\hline
\textbf{Technology} & \textbf{Version} & \textbf{Purpose} \\
\hline
Prometheus Client & $\geq$ 0.20.0 & Metrics collection and exposition \\
structlog & $\geq$ 24.1.0 & Structured JSON logging \\
OpenTelemetry API & $\geq$ 1.26.0 & Distributed tracing \\
OpenTelemetry SDK & $\geq$ 1.26.0 & Trace sampling and export \\
OTLP Exporter & $\geq$ 1.26.0 & Trace export to collectors \\
Grafana & 11.0.0 & Metrics visualization and dashboards \\
\hline
\end{tabular}
\end{table}

\subsection{LLM Provider Integration}

The platform supports multiple LLM providers through a unified adapter layer:

\begin{table}[ht]
\centering
\caption{Supported LLM providers and configurations.}
\label{tab:llm-providers}
\small
\begin{tabular}{lll}
\hline
\textbf{Provider} & \textbf{Protocol} & \textbf{Use Case} \\
\hline
Ollama & OpenAI-compatible API & Local/self-hosted models \\
OpenAI & OpenAI API & Cloud-hosted GPT models \\
Azure OpenAI & Azure API & Enterprise cloud deployment \\
Custom & Configurable & Organization-specific endpoints \\
\hline
\end{tabular}
\end{table}

Ollama is the primary provider for local deployment, with extensive configuration options:

\begin{lstlisting}[language=Python, caption={Ollama integration settings.}]
OLLAMA_BASE_URL: str | None = Field(default=None)
OLLAMA_TIMEOUT_SECS: int = Field(default=180)
LLM_WARMUP_TIMEOUT: int = Field(default=300)
DEFAULT_MODEL_NAME: str = Field(default="phi3:mini")
LLM_DEVICE: str = Field(default="cpu")
LLM_MAX_TOKENS: int = Field(default=2048)
LLM_MAX_STEPS: int = Field(default=10)
\end{lstlisting}

\subsection{Presentation Layer Technologies}

The platform provides two distinct user interfaces optimized for different personas.

\subsubsection{Agent Chat UI (Next.js)}

The conversational interface for end-users is built with modern React technologies:

\begin{table}[ht]
\centering
\caption{Agent Chat UI technology stack.}
\label{tab:ui-agent-stack}
\small
\begin{tabular}{lll}
\hline
\textbf{Technology} & \textbf{Version} & \textbf{Purpose} \\
\hline
Next.js & 14.2.15 & React framework with App Router \\
React & 18.3.1 & UI component library \\
TypeScript & 5.6.3 & Type-safe JavaScript \\
Zustand & 4.5.5 & Lightweight state management \\
Tailwind CSS & 3.4.14 & Utility-first styling \\
Radix UI & Latest & Accessible UI primitives \\
Lucide React & 0.451.0 & Icon library \\
\hline
\end{tabular}
\end{table}

\subsubsection{Control Panel UI (Streamlit)}

The administrative interface leverages Python-native web application development:

\begin{table}[ht]
\centering
\caption{Control Panel UI technology stack.}
\label{tab:ui-control-stack}
\small
\begin{tabular}{lll}
\hline
\textbf{Technology} & \textbf{Version} & \textbf{Purpose} \\
\hline
Streamlit & $\geq$ 1.30.0 & Python web application framework \\
Pandas & $\geq$ 2.2.0 & Data manipulation and display \\
httpx & $\geq$ 0.27.0 & HTTP client for API communication \\
\hline
\end{tabular}
\end{table}

\subsection{Background Processing}

Asynchronous job processing is implemented through:

\begin{itemize}
    \item \textbf{APScheduler} ($\geq$ 3.10.4): Background task scheduling with cron-style triggers
    \item \textbf{Redis queues}: BRPOP-based job queue with priority support
    \item \textbf{Custom worker process}: Dedicated container for job execution
\end{itemize}

\subsection{Container Orchestration}

The platform is containerized using Docker with orchestration through Docker Compose:

\begin{table}[ht]
\centering
\caption{Docker Compose service definitions.}
\label{tab:docker-services}
\small
\begin{tabular}{llp{6cm}}
\hline
\textbf{Service} & \textbf{Image} & \textbf{Purpose} \\
\hline
app & Custom (Dockerfile) & FastAPI backend application \\
worker & Custom (Dockerfile) & Background job processor \\
postgres & postgres:16-alpine & Control plane database \\
redis & redis:7-alpine & Cache and queue storage \\
memgraph & memgraph-platform:latest & Graph database \\
ollama & ollama/ollama:latest & Local LLM inference \\
prometheus & prom/prometheus:latest & Metrics collection \\
grafana & grafana/grafana:11.0.0 & Metrics visualization \\
ui\_control\_panel & Custom (Dockerfile) & Streamlit admin interface \\
\hline
\end{tabular}
\end{table}

\subsection{Development and Quality Tools}

The development toolchain ensures code quality and consistency:

\begin{table}[ht]
\centering
\caption{Development and quality assurance tools.}
\label{tab:dev-tools}
\small
\begin{tabular}{lll}
\hline
\textbf{Tool} & \textbf{Version} & \textbf{Purpose} \\
\hline
Ruff & $\geq$ 0.5.0 & Linting and import sorting \\
Black & $\geq$ 24.4.2 & Code formatting \\
mypy & $\geq$ 1.10.0 & Static type checking \\
Bandit & $\geq$ 1.7.9 & Security vulnerability scanning \\
pytest & $\geq$ 8.2.0 & Test framework \\
coverage & $\geq$ 7.5.0 & Code coverage measurement \\
\hline
\end{tabular}
\end{table}

The \texttt{pyproject.toml} configuration centralizes all tool settings, enabling consistent enforcement across development environments:

\begin{lstlisting}[language=Python, caption={Code quality configuration excerpt from \texttt{pyproject.toml}.}]
[tool.ruff]
target-version = "py311"
line-length = 100

[tool.ruff.lint]
select = [
    "E", "F", "W",   # Core style checks
    "I",             # isort import ordering
    "B",             # bugbear security patterns
    "SIM",           # simplification suggestions
    "RUF",           # ruff-specific rules
]

[tool.mypy]
python_version = "3.11"
plugins = ["pydantic.mypy"]
check_untyped_defs = true
disallow_untyped_defs = true
\end{lstlisting}

